\chapter{Inversión}

\titular{La inversión física propuesta es de \textbf{2,5\,\% del PIB} y la financiera de
\textbf{0,7}, el nivel más alto de la serie proyectada. La segunda vuelve a 0,3 en 2027.
\textbf{Dentro de los capítulos 6000 y 7000 del objeto del gasto, cerca de la mitad
corresponde a inversiones financieras y provisiones, con la aportación de capital a Pemex como
componente dominante, y esa aportación no se repite.}}

\quecambio{La obra pública directa ---capítulo 6000 del objeto del gasto--- suma
\textbf{431\,608,6 millones de pesos}, 1,11\,\% del PIB. Las inversiones financieras y otras
provisiones ---capítulo 7000--- suman \textbf{388\,161,3}, y dentro de ellas
\textbf{266\,053,3 son compra de títulos y valores}.}

\porqueimporta{Inversión física y gasto de inversión son cosas distintas y el paquete las
presenta juntas. \textbf{De cada peso de los capítulos 6000 y 7000 en 2026, cerca de
cuarenta y siete centavos son inversiones financieras y provisiones}, no obra. Un lector que
tome la suma de esos dos capítulos como capacidad productiva instalada se equivoca por casi la
mitad.}

\section{Declaración de perímetro}

\marginal{obra y títulos son cosas distintas}
Este capítulo separa dos objetos que la clasificación económica agrupa:

\begin{list}{}{\setlength{\leftmargin}{1.4em}\setlength{\itemsep}{2pt}}
\item[$\bullet$] \textbf{Inversión pública (capítulo 6000):} obra pública en bienes propios y
de dominio público. 431\,608,6 millones.
\item[$\bullet$] \textbf{Inversiones financieras y otras provisiones (capítulo 7000):}
266\,053,3 de compra de títulos y valores más 122\,108,0 de provisiones para contingencias.
Total 388\,161,3 millones.
\end{list}

Los dos son gasto de capital y \textbf{solo el primero construye algo}.

\section{La compra de títulos}

De los 266\,053,3 millones de compra de títulos y valores, \textbf{263\,476,3 están en el Ramo
18, Energía}, en una sola partida específica. Es la aportación de capital a Petróleos
Mexicanos que el capítulo 3 describe.

\textbf{No es un juicio decir que no es obra: es una clasificación del propio decreto.} Y el
marco de mediano plazo lo confirma como movimiento de un solo año: la inversión financiera
pasa de 0,4\,\% del PIB en 2025 a \textbf{0,7 en 2026} y vuelve a \textbf{0,3 en 2027}, donde
se queda hasta 2031.

\section{La obra, y lo que se puede decir de ella}

Los 431\,608,6 millones de obra pública se pueden seguir por su clave de cartera de inversión,
que la base del proyecto publica. Los proyectos identificados mayores valen 30\,000,0,
27\,720,0, 15\,875,7 y 14\,386,0 millones de pesos.

\textbf{Aquí hay una limitación que hay que declarar.} La base trae la clave de cartera y
\textbf{no su denominación}. Se pueden publicar los montos por clave y no se puede decir a qué
proyecto corresponde cada uno sin el tomo de proyectos de inversión, que no está en carpeta.
Este documento publica las claves y sus montos, y no les pone nombre.

Además, \textbf{11\,260\,234,1 millones del gasto bruto tienen clave de cartera cero}, es
decir, no corresponden a un proyecto de inversión registrado. Es la inmensa mayoría del
presupuesto y es normal ---el gasto corriente no tiene cartera---, pero conviene decirlo para
que nadie lea la cartera como una partición del gasto.

\section{La obra que ejecutan las fuerzas armadas}

El Ramo 07 mantiene \textbf{31\,620,0 millones de pesos en la función 3.5, Transporte},
después de 41\,772,1 en el proyecto anterior: una caída de 27,7\,\% real. Y el Grupo
Aeroportuario y Ferroviario que administra sube de 2\,275,5 a 3\,441,9 millones.

La obra ferroviaria en conjunto ---la que vive en el Ramo 09--- pasa de 94\,352,0 a
\textbf{95\,966,5 millones}: \textbf{4,0\,\% menos en términos reales}, con un cambio de
unidad responsable de por medio, el tercero en tres ejercicios. \textbf{El total apenas se
mueve; lo que se mueve es quién lo ejecuta.}

\section{Implicaciones}

La inversión física de 2,5\,\% del PIB es el punto más alto del horizonte proyectado y desde
ahí baja: 2,6 en 2027 y 2,3 de 2029 en adelante. \textbf{El escenario oficial de estabilidad
de la deuda descansa, entre otras cosas, sobre una inversión pública que no vuelve a los
niveles actuales en seis años.}

Y la suma de los capítulos 6000 y 7000 de 2026 no es lo que la palabra «capital» sugiere:
cerca de la mitad son inversiones financieras y provisiones, dominadas por una sola operación
con una empresa del Estado. \textbf{Un documento que reporte «la inversión sube» sin
separar los dos objetos estará describiendo una recapitalización como si fuera infraestructura.}

\begin{ficha}{inversión}
\fila{Perímetro}{Capítulos 6000 y 7000 del objeto del gasto, gasto bruto de Gobierno Federal
más entidades, presentados por separado y nunca sumados sin etiqueta.}
\fila{Línea de comparación}{\textbf{Línea P} para las cifras por ramo y función. Los
agregados por capítulo son de 2026 y no tienen contraparte de proyecto 2025 en la base
publicada, y se presentan como nivel.}
\fila{Deflactor}{Del PIB, 4,8\,\%, CGPE 2026.}
\fila{Año de los pesos}{Corrientes.}
\fila{PIB y añada}{38\,715,9 miles de millones, CGPE 2026.}
\fila{Denominador demográfico}{No aplica.}
\fila{Fuentes}{Base del proyecto 2026 por capítulo, concepto, partida específica y clave de
cartera; analíticos del proyecto por función; CGPE 2026 Anexo III.2.}
\fila{Qué no puede afirmarse}{El nombre de los proyectos de la cartera: la base publica la
clave y no la denominación, y el tomo correspondiente no está en carpeta. La composición de la
aportación de capital a Pemex entre amortización e inversión productiva. Y el avance físico de
ningún proyecto, que no es materia de un paquete.}
\end{ficha}
\clearpage
